\pagestyle{empty}
\tight
\begin{tblr}{width=\textwidth,colspec={|Q[c,m,1.9cm]|X[l,m]|},hlines,vlines,hspan=minimal,row{1}={bg=headblue},rowsep=1pt,colsep=3pt}
\SetCell{c,m}\hfil\logo\hfil & \textbf{POLITEKNIK NEGERI MALANG}\par\textbf{JURUSAN TEKNOLOGI INFORMASI}\par\textbf{PROGRAM STUDI: D4 TEKNIK INFORMATIKA} \\
\SetCell[c=2]{c} \textbf{RENCANA PEMBELAJARAN SEMESTER (RPS)} & \\
\end{tblr}
\begin{longtblr}{width=\textwidth,colspec={|Q[l,m,1.9cm]|X[0.85,l,m]|X[1.45,l,m]|X[0.75,l,m]|X[1.15,l,m]|},hlines,vlines,hspan=minimal,rowsep=1pt,colsep=2pt,row{1,3}={bg=headgray,font=\bfseries}}
MATA KULIAH & KODE & BOBOT (SKS)/jam & SEMESTER & TGL. PENYUSUNAN \\
Kewirausahaan Berbasis Teknologi & RTI255001 & 2 SKS/4 jam & 5 & 6 Januari 2026 \\
OTORISASI & Dosen Pengembang RPS & Koordinator RMK & \SetCell[c=2]{l} Ka PRODI &  \\
 & Tim Dosen Mata Kuliah &  & \SetCell[c=2]{l}{\vspace{8mm}\par Dr. Ely Setyo Astuti, S.T., M.T.} &  \\
\SetCell[r=11]{l} Capaian Pembelajaran (CP) & \SetCell[c=4]{l,bg=headgray,font=\bfseries} Capaian Pembelajaran Lulusan Program Studi (CPL-Prodi) &  &  &  \\
 & CPL01 & \SetCell[c=3]{l} Mampu menerapkan nilai ketakwaan, kepatuhan hukum, disiplin, profesionalisme, etika profesi, dan kewirausahaan dalam praktik profesi teknologi informasi. &  &  \\
 & CPL03 & \SetCell[c=3]{l} Mampu mengelola diri, bekerja profesional, berkomunikasi lisan dan tulisan, serta melakukan presentasi secara efektif. &  &  \\
 & \SetCell[c=4]{l,bg=headgray,font=\bfseries} Capaian Pembelajaran mata kuliah (CPL-MK) &  &  &  \\
 & CPMK0101 & \SetCell[c=3]{l} Mahasiswa mampu menerapkan etika profesi, jiwa kewirausahaan berbasis teknologi, dan tanggung jawab sosial dalam pengembangan produk dan bisnis TI. &  &  \\
 & CPMK0301 & \SetCell[c=3]{l} Mahasiswa mampu mengelola diri secara profesional, berkomunikasi efektif, dan mempresentasikan rencana bisnis teknologi kepada pemangku kepentingan. &  &  \\
 & \SetCell[c=4]{l,bg=headgray,font=\bfseries} Kemampuan akhir tiap tahapan belajar (Sub-CPMK) &  &  &  \\
 & SCPMK0101-03701 & \SetCell[c=3]{l} Mahasiswa mampu menjelaskan konsep kewirausahaan, mengidentifikasi karakteristik wirausaha, mengembangkan intensi kewirausahaan, serta membangkitkan dan memilih ide produk teknologi. &  &  \\
 & SCPMK0101-03702 & \SetCell[c=3]{l} Mahasiswa mampu menganalisis permintaan pasar, menyusun rencana pemasaran, mengidentifikasi aspek teknis bisnis, serta menyusun laporan keuangan dasar dan penetapan harga. &  &  \\
 & SCPMK0301-03703 & \SetCell[c=3]{l} Mahasiswa mampu merancang Business Model Canvas, menganalisis sumber pendanaan, aspek hukum, tanggung jawab sosial, manajemen SDM, dan strategi pengembangan produk baru. &  &  \\
 & SCPMK0301-03704 & \SetCell[c=3]{l} Mahasiswa mampu menyusun dan mempresentasikan rencana bisnis teknologi secara profesional dengan argumentasi berbasis data kepada pemangku kepentingan. &  &  \\
\SetCell[r=2]{l} Penilaian & \SetCell[c=4]{l} *Nilai CPMK didapat dari agregasi nilai SUB-CPMK yang dibebankan &  &  &  \\
 & \SetCell[c=4]{l}{\tiny\begin{tblr}{width=\linewidth,colspec={|Q[c,m,0.1\linewidth]|Q[c,m,0.16\linewidth]|X[l,m]|X[l,m]|X[l,m]|X[l,m]|X[l,m]|Q[c,m,0.08\linewidth]|},hlines,vlines,hspan=minimal,rowsep=1pt,colsep=0.7pt,row{1,2}={bg=headgray,font=\bfseries}}
Kode CPMK & Kode SCPMK & \SetCell[c=5]{c} Bobot per Bentuk Penilaian & & & & & Total Bobot \\
 & & CM & Tugas & Kuis & UTS & UAS & \\
\SetCell[r=2]{c} CPMK0101 & SCPMK0101-03701 & 13\%: Business Game, partisipasi, PEC & 5\%: ide produk dan profil wirausaha & 5\%: konsep kewirausahaan & 5\%: kewirausahaan dan ide bisnis & -- & 28\% \\
 & SCPMK0101-03702 & -- & 23\%: analisis pasar, aspek teknis, keuangan & 5\%: pasar dan teknis bisnis & 5\%: strategi pemasaran dan keuangan & -- & 33\% \\
\SetCell[r=2]{c} CPMK0301 & SCPMK0301-03703 & 12\%: workshop BMC, pendanaan, hukum & 7\%: BMC, pendanaan, SDM, produk baru & -- & -- & -- & 19\% \\
 & SCPMK0301-03704 & -- & -- & -- & -- & 20\%: presentasi rencana bisnis & 20\% \\
\SetCell[c=7]{r}\textbf{Total Per Penilaian} & & & & & & & \textbf{100\%} \\
\end{tblr}} &  &  &  \\
Diskripsi Singkat Mata Kuliah & \SetCell[c=4]{l} Mata kuliah Kewirausahaan Berbasis Teknologi membangun jiwa dan kompetensi kewirausahaan mahasiswa melalui pemahaman konsep kewirausahaan, simulasi bisnis (Business Game), analisis pasar, perencanaan keuangan, penyusunan Business Model Canvas, serta pengembangan dan presentasi rencana bisnis berbasis teknologi informasi. &  &  &  \\
Bahan Kajian / Materi Pembelajaran & \SetCell[c=4]{l}\begin{enumerate}\item Konsep kewirausahaan, karakteristik wirausaha, dan intensi kewirausahaan\item Analisis pasar, aspek teknis bisnis, dan manajemen keuangan dasar\item Business Model Canvas, sumber pendanaan, dan aspek legal bisnis\item Manajemen SDM dan pengembangan produk baru\item Penyusunan dan presentasi rencana bisnis teknologi\end{enumerate} &  &  &  \\
\SetCell[r=2]{l} Pustaka & \SetCell[c=4]{l} \textbf{Utama:}\begin{enumerate}\item Hisrich, R.D., Peters, M.P., Shepherd, D.A. 2017. Entrepreneurship. 10th ed. McGraw-Hill Education.\item Osterwalder, A., Pigneur, Y. 2010. Business Model Generation. John Wiley \& Sons.\item Bygrave, W., Zacharakis, A. 2011. Entrepreneurship. 2nd ed. John Wiley \& Sons.\end{enumerate} &  &  &  \\
 & \SetCell[c=4]{l} \textbf{Pendukung:} Materi LMS, modul business game, jurnal kewirausahaan, dan referensi dosen. &  &  &  \\
Media Pembelajaran & \SetCell[c=2]{l} \textbf{Software:} Microsoft Office (Word, Excel, PowerPoint), Canva/Figma (opsional), LMS. &  & \SetCell[c=2]{l} \textbf{Hardware:} Komputer/laptop, proyektor. &  \\
Nama Dosen Pengampu & \SetCell[c=4]{l} Tim Dosen Mata Kuliah &  &  &  \\
Matakuliah Syarat & \SetCell[c=4]{l} - &  &  &  \\
\end{longtblr}
\begin{longtblr}{width=\textwidth,colspec={|Q[c,m,0.06\textwidth]|X[1.35,l,m]|X[1.08,l,m]|X[1.16,l,m]|X[0.7,l,m]|X[1.24,l,m]|X[1.06,l,m]|X[0.98,l,m]|Q[c,m,0.05\textwidth]|},hlines,vlines,hspan=minimal,rowhead=2,row{1,2}={bg=headgreen,font=\bfseries},rowsep=1pt,colsep=1.5pt}
Mgg.\par Ke & Kemampuan Akhir Yang Direncanakan (Sub-CP-MK) & Bahan kajian (Materi Pembelajaran) & Bentuk dan Metode Pembelajaran & Estimasi Waktu & Pengalaman Belajar Mahasiswa & Kriteria \& Bentuk Penilaian & Indikator Penilaian & Bobot Penilaian (\%) \\
(1) & (2) & (3) & (4) & (5) & (6) & (7) & (8) & (9) \\
1 & SCPMK0101-03701: Mahasiswa mampu menjelaskan konsep kewirausahaan, mengidentifikasi karakteristik wirausaha, mengembangkan intensi kewirausahaan, serta membangkitkan dan memilih ide produk teknologi. & Usaha, kewirausahaan, intensi kewirausahaan dan wirausaha. & Modalitas: Blended Learning. Bentuk: Luring/Daring. Metode: Blended Learning, luring/daring, experiential learning, business simulation, studi kasus, presentasi. & 1 x 2 x 50' tatap muka; 1 x 2 x 50' tugas/praktik mandiri. & Mahasiswa mengerjakan aktivitas terarah untuk usaha, kewirausahaan, intensi kewirausahaan dan wirausaha dan menunjukkan evidence hasil belajar. & Bentuk penilaian: CM (partisipasi). Mengacu rubrik RTM. & Ketepatan konsep; kualitas artefak/praktik; validasi dan dokumentasi. & 3 \\
2 & SCPMK0101-03701: Mahasiswa mampu menjelaskan konsep kewirausahaan, mengidentifikasi karakteristik wirausaha, mengembangkan intensi kewirausahaan, serta membangkitkan dan memilih ide produk teknologi. & Karakteristik wirausaha; Business Game Module 1. & Modalitas: Blended Learning. Bentuk: Luring/Daring. Metode: Blended Learning, luring/daring, experiential learning, business simulation, studi kasus, presentasi. & 1 x 2 x 50' tatap muka; 1 x 2 x 50' tugas/praktik mandiri. & Mahasiswa mengerjakan aktivitas terarah untuk karakteristik wirausaha dan business game module 1 dan menunjukkan evidence hasil belajar. & Bentuk penilaian: CM (Business Game). Mengacu rubrik RTM. & Ketepatan konsep; kualitas artefak/praktik; validasi dan dokumentasi. & 5 \\
3 & SCPMK0101-03701: Mahasiswa mampu menjelaskan konsep kewirausahaan, mengidentifikasi karakteristik wirausaha, mengembangkan intensi kewirausahaan, serta membangkitkan dan memilih ide produk teknologi. & Memperbaiki Personal Entrepreneurial Characteristics (PEC); Business Game Module 2. & Modalitas: Blended Learning. Bentuk: Luring/Daring. Metode: Blended Learning, luring/daring, experiential learning, business simulation, studi kasus, presentasi. & 1 x 2 x 50' tatap muka; 1 x 2 x 50' tugas/praktik mandiri. & Mahasiswa mengerjakan aktivitas terarah untuk memperbaiki PEC dan business game module 2 dan menunjukkan evidence hasil belajar. & Bentuk penilaian: CM (Business Game). Mengacu rubrik RTM. & Ketepatan konsep; kualitas artefak/praktik; validasi dan dokumentasi. & 5 \\
4 & SCPMK0101-03701: Mahasiswa mampu menjelaskan konsep kewirausahaan, mengidentifikasi karakteristik wirausaha, mengembangkan intensi kewirausahaan, serta membangkitkan dan memilih ide produk teknologi. & Membangkitkan dan memilih ide produk dan ide bisnis. & Modalitas: Blended Learning. Bentuk: Luring/Daring. Metode: Blended Learning, luring/daring, experiential learning, business simulation, studi kasus, presentasi. & 1 x 2 x 50' tatap muka; 1 x 2 x 50' tugas/praktik mandiri. & Mahasiswa mengerjakan aktivitas terarah untuk membangkitkan dan memilih ide produk dan ide bisnis dan menunjukkan evidence hasil belajar. & Bentuk penilaian: Tugas (ide produk dan profil wirausaha). Mengacu rubrik RTM. & Ketepatan konsep; kualitas artefak/praktik; validasi dan dokumentasi. & 5 \\
5 & SCPMK0101-03702: Mahasiswa mampu menganalisis permintaan pasar, menyusun rencana pemasaran, mengidentifikasi aspek teknis bisnis, serta menyusun laporan keuangan dasar dan penetapan harga. & Memahami pasar, mengukur permintaan pasar, dan rencana pemasaran. & Modalitas: Blended Learning. Bentuk: Luring/Daring. Metode: Blended Learning, luring/daring, experiential learning, business simulation, studi kasus, presentasi. & 1 x 2 x 50' tatap muka; 1 x 2 x 50' tugas/praktik mandiri. & Mahasiswa mengerjakan aktivitas terarah untuk memahami pasar, mengukur permintaan pasar, dan rencana pemasaran dan menunjukkan evidence hasil belajar. & Bentuk penilaian: Tugas (analisis pasar). Mengacu rubrik RTM. & Ketepatan konsep; kualitas artefak/praktik; validasi dan dokumentasi. & 5 \\
6 & SCPMK0101-03702: Mahasiswa mampu menganalisis permintaan pasar, menyusun rencana pemasaran, mengidentifikasi aspek teknis bisnis, serta menyusun laporan keuangan dasar dan penetapan harga. & Aspek-aspek teknis bisnis. & Modalitas: Blended Learning. Bentuk: Luring/Daring. Metode: Blended Learning, luring/daring, experiential learning, business simulation, studi kasus, presentasi. & 1 x 2 x 50' tatap muka; 1 x 2 x 50' tugas/praktik mandiri. & Mahasiswa mengerjakan aktivitas terarah untuk aspek-aspek teknis bisnis dan menunjukkan evidence hasil belajar. & Bentuk penilaian: Tugas (aspek teknis bisnis). Mengacu rubrik RTM. & Ketepatan konsep; kualitas artefak/praktik; validasi dan dokumentasi. & 5 \\
7 & SCPMK0101-03701, SCPMK0101-03702 & Kuis; Memahami dan membuat buku kas. & Modalitas: Blended Learning. Bentuk: Luring/Daring. Metode: Blended Learning, luring/daring, experiential learning, business simulation, studi kasus, presentasi. & 1 x 2 x 50' tatap muka; 1 x 2 x 50' tugas/praktik mandiri. & Mahasiswa mengerjakan aktivitas terarah untuk kuis dan memahami serta membuat buku kas dan menunjukkan evidence hasil belajar. & Bentuk penilaian: Kuis (konsep kewirausahaan, pasar, teknis). Mengacu rubrik RTM. & Ketepatan konsep; kualitas artefak/praktik; validasi dan dokumentasi. & 10 \\
8 & SCPMK0101-03701, SCPMK0101-03702 & UTS: kewirausahaan, analisis pasar, strategi pemasaran, dan keuangan dasar bisnis. & Modalitas: Blended Learning. Bentuk: Luring/Daring. Metode: Blended Learning, luring/daring, experiential learning, business simulation, studi kasus, presentasi. & 1 x 2 x 50' tatap muka; 1 x 2 x 50' tugas/praktik mandiri. & Mahasiswa mengerjakan aktivitas terarah untuk uts dan menunjukkan evidence hasil belajar. & Bentuk penilaian: UTS (aspek teknis dan keuangan bisnis). Mengacu rubrik RTM. & Ketepatan konsep; kualitas artefak/praktik; validasi dan dokumentasi. & 10 \\
9 & SCPMK0101-03702: Mahasiswa mampu menganalisis permintaan pasar, menyusun rencana pemasaran, mengidentifikasi aspek teknis bisnis, serta menyusun laporan keuangan dasar dan penetapan harga. & Menghitung biaya, menetapkan harga jual, dan menyusun laba rugi. & Modalitas: Blended Learning. Bentuk: Luring/Daring. Metode: Blended Learning, luring/daring, experiential learning, business simulation, studi kasus, presentasi. & 1 x 2 x 50' tatap muka; 1 x 2 x 50' tugas/praktik mandiri. & Mahasiswa mengerjakan aktivitas terarah untuk menghitung biaya, menetapkan harga jual, dan menyusun laba rugi dan menunjukkan evidence hasil belajar. & Bentuk penilaian: Tugas (keuangan: biaya dan harga). Mengacu rubrik RTM. & Ketepatan konsep; kualitas artefak/praktik; validasi dan dokumentasi. & 5 \\
10 & SCPMK0101-03702: Mahasiswa mampu menganalisis permintaan pasar, menyusun rencana pemasaran, mengidentifikasi aspek teknis bisnis, serta menyusun laporan keuangan dasar dan penetapan harga. & Modal awal, proyeksi laporan keuangan, dan ratio keuangan dasar (1). & Modalitas: Blended Learning. Bentuk: Luring/Daring. Metode: Blended Learning, luring/daring, experiential learning, business simulation, studi kasus, presentasi. & 1 x 2 x 50' tatap muka; 1 x 2 x 50' tugas/praktik mandiri. & Mahasiswa mengerjakan aktivitas terarah untuk modal awal, proyeksi laporan keuangan, dan ratio keuangan dasar bagian 1 dan menunjukkan evidence hasil belajar. & Bentuk penilaian: Tugas (proyeksi keuangan 1). Mengacu rubrik RTM. & Ketepatan konsep; kualitas artefak/praktik; validasi dan dokumentasi. & 4 \\
11 & SCPMK0101-03702: Mahasiswa mampu menganalisis permintaan pasar, menyusun rencana pemasaran, mengidentifikasi aspek teknis bisnis, serta menyusun laporan keuangan dasar dan penetapan harga. & Modal awal, proyeksi laporan keuangan, dan ratio keuangan dasar (2). & Modalitas: Blended Learning. Bentuk: Luring/Daring. Metode: Blended Learning, luring/daring, experiential learning, business simulation, studi kasus, presentasi. & 1 x 2 x 50' tatap muka; 1 x 2 x 50' tugas/praktik mandiri. & Mahasiswa mengerjakan aktivitas terarah untuk modal awal, proyeksi laporan keuangan, dan ratio keuangan dasar bagian 2 dan menunjukkan evidence hasil belajar. & Bentuk penilaian: Tugas (proyeksi keuangan 2). Mengacu rubrik RTM. & Ketepatan konsep; kualitas artefak/praktik; validasi dan dokumentasi. & 4 \\
12 & SCPMK0301-03703: Mahasiswa mampu merancang Business Model Canvas, menganalisis sumber pendanaan, aspek hukum, tanggung jawab sosial, manajemen SDM, dan strategi pengembangan produk baru. & Business Model Canvas (1). & Modalitas: Blended Learning. Bentuk: Luring/Daring. Metode: Blended Learning, luring/daring, experiential learning, business simulation, studi kasus, presentasi. & 1 x 2 x 50' tatap muka; 1 x 2 x 50' tugas/praktik mandiri. & Mahasiswa mengerjakan aktivitas terarah untuk business model canvas bagian 1 dan menunjukkan evidence hasil belajar. & Bentuk penilaian: CM (workshop BMC). Mengacu rubrik RTM. & Ketepatan konsep; kualitas artefak/praktik; validasi dan dokumentasi. & 4 \\
13 & SCPMK0301-03703: Mahasiswa mampu merancang Business Model Canvas, menganalisis sumber pendanaan, aspek hukum, tanggung jawab sosial, manajemen SDM, dan strategi pengembangan produk baru. & Business Model Canvas (2). & Modalitas: Blended Learning. Bentuk: Luring/Daring. Metode: Blended Learning, luring/daring, experiential learning, business simulation, studi kasus, presentasi. & 1 x 2 x 50' tatap muka; 1 x 2 x 50' tugas/praktik mandiri. & Mahasiswa mengerjakan aktivitas terarah untuk business model canvas bagian 2 dan menunjukkan evidence hasil belajar. & Bentuk penilaian: CM (workshop BMC). Mengacu rubrik RTM. & Ketepatan konsep; kualitas artefak/praktik; validasi dan dokumentasi. & 4 \\
14 & SCPMK0301-03703: Mahasiswa mampu merancang Business Model Canvas, menganalisis sumber pendanaan, aspek hukum, tanggung jawab sosial, manajemen SDM, dan strategi pengembangan produk baru. & Sumber-sumber pendanaan untuk bisnis baru, bank, dan kebutuhan pendanaan. & Modalitas: Blended Learning. Bentuk: Luring/Daring. Metode: Blended Learning, luring/daring, experiential learning, business simulation, studi kasus, presentasi. & 1 x 2 x 50' tatap muka; 1 x 2 x 50' tugas/praktik mandiri. & Mahasiswa mengerjakan aktivitas terarah untuk sumber-sumber pendanaan untuk bisnis baru dan menunjukkan evidence hasil belajar. & Bentuk penilaian: Tugas (pendanaan dan SDM). Mengacu rubrik RTM. & Ketepatan konsep; kualitas artefak/praktik; validasi dan dokumentasi. & 3 \\
15 & SCPMK0301-03703: Mahasiswa mampu merancang Business Model Canvas, menganalisis sumber pendanaan, aspek hukum, tanggung jawab sosial, manajemen SDM, dan strategi pengembangan produk baru. & Aspek hukum dan tanggung jawab sosial. & Modalitas: Blended Learning. Bentuk: Luring/Daring. Metode: Blended Learning, luring/daring, experiential learning, business simulation, studi kasus, presentasi. & 1 x 2 x 50' tatap muka; 1 x 2 x 50' tugas/praktik mandiri. & Mahasiswa mengerjakan aktivitas terarah untuk aspek hukum dan tanggung jawab sosial dan menunjukkan evidence hasil belajar. & Bentuk penilaian: CM (aspek hukum). Mengacu rubrik RTM. & Ketepatan konsep; kualitas artefak/praktik; validasi dan dokumentasi. & 4 \\
16 & SCPMK0301-03703: Mahasiswa mampu merancang Business Model Canvas, menganalisis sumber pendanaan, aspek hukum, tanggung jawab sosial, manajemen SDM, dan strategi pengembangan produk baru. & Manajemen sumber daya manusia, jaringan, dan mengembangkan produk baru. & Modalitas: Blended Learning. Bentuk: Luring/Daring. Metode: Blended Learning, luring/daring, experiential learning, business simulation, studi kasus, presentasi. & 1 x 2 x 50' tatap muka; 1 x 2 x 50' tugas/praktik mandiri. & Mahasiswa mengerjakan aktivitas terarah untuk manajemen SDM, jaringan, dan pengembangan produk baru dan menunjukkan evidence hasil belajar. & Bentuk penilaian: Tugas (SDM dan produk baru). Mengacu rubrik RTM. & Ketepatan konsep; kualitas artefak/praktik; validasi dan dokumentasi. & 4 \\
17 & SCPMK0301-03704: Mahasiswa mampu menyusun dan mempresentasikan rencana bisnis teknologi secara profesional dengan argumentasi berbasis data kepada pemangku kepentingan. & UAS: Presentasi Rencana Bisnis Teknologi. & Modalitas: Blended Learning. Bentuk: Luring/Daring. Metode: Blended Learning, luring/daring, experiential learning, business simulation, studi kasus, presentasi. & 1 x 2 x 50' tatap muka; 1 x 2 x 50' tugas/praktik mandiri. & Mahasiswa mengerjakan aktivitas terarah untuk presentasi rencana bisnis teknologi dan menunjukkan evidence hasil belajar. & Bentuk penilaian: UAS (presentasi rencana bisnis teknologi). Mengacu rubrik RTM. & Ketepatan konsep; kualitas artefak/praktik; validasi dan dokumentasi. & 20 \\
\end{longtblr}
